\section{候选方向}
% 8
\begin{frame}[t]{两个候选方向：共用一套数据与评分流程}
\vspace{.1cm}
\centering
\begin{tikzpicture}[card/.style={box,align=left,text width=6.2cm,inner sep=9pt}]
 \node[card,fill=PaleTeal,draw=DarkBlue,line width=.9pt] (A) at (0,0) {%
  \ico{\faStar}\enspace\textbf{方向 A\quad 跨背景扰动响应预测}\\[.12cm]
  {\small\color{DarkBlue}只有对照数据的新细胞背景中，预训练细胞模型能否超过简单基线？}\\[.14cm]
  \kv{数据}{多背景 Perturb-seq、VCC 2025 数据}\\[.04cm]
  \kv{交付}{基线表、跨背景增益与分层分析、可复现脚本}\\[.04cm]
  \kv{医院衔接}{患者来源样本作为新的“目标背景”}};
 \node[card] (B) at (7.2,0) {%
  \ico{\faSearch}\enspace\textbf{方向 B\quad 扰动响应的可预测性诊断}\\[.12cm]
  {\small\color{DarkBlue}扰动响应中哪些部分可预测？把系统性变化与扰动特异信号分开评估。}\\[.14cm]
  \kv{数据}{与 A 相同，加 Systema 评估框架}\\[.04cm]
  \kv{交付}{扰动 × 基因模块的可预测性地图、可信预测清单}\\[.04cm]
  \kv{医院衔接}{告诉合作者哪些预测值得在样本上验证}};
 \node[fill=DarkBlue,text=white,font=\scriptsize\bfseries,rounded corners=2pt,inner sep=3pt,anchor=east] at ([xshift=-.3cm]A.north east) {首选};
 \node[fill=Muted,text=white,font=\scriptsize\bfseries,rounded corners=2pt,inner sep=3pt,anchor=east] at ([xshift=-.3cm]B.north east) {备选 · 可并行};
 \node[font=\scriptsize\bfseries,text=Muted,anchor=west] (sl) at (-3.4,-2.75) {\faCogs\enspace 共用流程};
 \node[chip,right=.25cm of sl] (f1) {数据审计};
 \node[chip,right=.25cm of f1] (f2) {按细胞背景划分};
 \node[chip,right=.25cm of f2] (f3) {评分脚本（冻结）};
 \node[chip,right=.25cm of f3] (f4) {全部基线};
 \node[chip,right=.25cm of f4] (f5) {分层分析};
 \draw[flow] (f1)--(f2);\draw[flow] (f2)--(f3);\draw[flow] (f3)--(f4);\draw[flow] (f4)--(f5);
\end{tikzpicture}
\remark{倾向先做 A；B 用到 A 的全部中间结果，可作为 A 的分析章节，也可独立成一个小题目。}
\source{两个方向均为候选研究设想，尚未实现；题目措辞待与老师讨论后确定。数据与协议来源见第 6 页。}
\end{frame}
